%=====================================================================
\chapter{Cloud and Edge Computing for Data Science}\label{ch:cloudedge}
%=====================================================================
\locator{5}{Part V --- Data Science in Systems Context}

\begin{outcomes}
By the end of this chapter you will be able to:
\begin{tight}
  \item \textbf{Distinguish} IaaS, PaaS and SaaS, and say who is responsible for what in each.
  \item \textbf{Explain} how containers differ from virtual machines and why that matters for reproducibility.
  \item \textbf{Choose} between cloud, fog and edge placement for a given inference workload.
  \item \textbf{Estimate} and control the cost of a data science workload.
  \item \textbf{Apply} the shared responsibility model to a cloud deployment.
\end{tight}
\end{outcomes}

\begin{prereq}
\chref{dataeng} (pipelines and storage). This chapter supplies the deployment
substrate for \chref{iot} and \chref{mlops}.
\end{prereq}

\begin{keyterms}
IaaS / PaaS / SaaS $\bullet$ public, private, hybrid cloud $\bullet$ elasticity
$\bullet$ virtual machine $\bullet$ container $\bullet$ image $\bullet$ Kubernetes
$\bullet$ serverless $\bullet$ GPU $\bullet$ cloud, fog and edge $\bullet$ latency
$\bullet$ quantisation $\bullet$ shared responsibility $\bullet$ data residency
$\bullet$ egress cost
\end{keyterms}

\section{Service Models}

\dsfig{%
\begin{tikzpicture}[font=\scriptsize,
  you/.style={rounded corners=2pt, fill=accent!22, draw=accent!60, text=accent!65!black,
              minimum width=2.5cm, minimum height=0.42cm, font=\tiny},
  them/.style={rounded corners=2pt, fill=secondary!16, draw=secondary!50, text=secondary!55!black,
               minimum width=2.5cm, minimum height=0.42cm, font=\tiny},
  hd/.style={rounded corners=3pt, fill=primary, text=white, font=\tiny\bfseries,
             minimum width=2.5cm, minimum height=0.5cm}]

\foreach \l/\i in {Application/0, Data/1, Runtime/2, Middleware/3, {OS}/4,
                   Virtualisation/5, Servers/6, Storage/7, Networking/8}{
  \pgfmathsetmacro{\yy}{2.15-\i*0.5}
  \node[font=\tiny, anchor=east, text=gray!50!black] at (-0.25,\yy) {\l};}

\node[hd] at (1.4,2.75)  {ON-PREMISES};
\node[hd] at (4.3,2.75)  {IaaS};
\node[hd] at (7.2,2.75)  {PaaS};
\node[hd] at (10.1,2.75) {SaaS};

% on-prem: all you
\foreach \i in {0,...,8}{\pgfmathsetmacro{\yy}{2.15-\i*0.5}\node[you] at (1.4,\yy) {};}
% IaaS: you down to OS
\foreach \i in {0,...,4}{\pgfmathsetmacro{\yy}{2.15-\i*0.5}\node[you] at (4.3,\yy) {};}
\foreach \i in {5,...,8}{\pgfmathsetmacro{\yy}{2.15-\i*0.5}\node[them] at (4.3,\yy) {};}
% PaaS: you app+data
\foreach \i in {0,1}{\pgfmathsetmacro{\yy}{2.15-\i*0.5}\node[you] at (7.2,\yy) {};}
\foreach \i in {2,...,8}{\pgfmathsetmacro{\yy}{2.15-\i*0.5}\node[them] at (7.2,\yy) {};}
% SaaS: you data only
\node[you] at (10.1,1.65) {};
\node[them] at (10.1,2.15) {};
\foreach \i in {2,...,8}{\pgfmathsetmacro{\yy}{2.15-\i*0.5}\node[them] at (10.1,\yy) {};}

\node[font=\scriptsize, anchor=west, align=left] at (11.6,1.9)
     {\textcolor{accent!65!black}{$\blacksquare$}~\textbf{you manage}\\[4pt]
      \textcolor{secondary!55!black}{$\blacksquare$}~provider manages};

\node[anchor=west, align=left, text width=6.6cm, font=\scriptsize, text=gray!55!black]
     at (-1.9,-2.5)
     {\textbf{The one row that never changes colour} is \emph{Data}. It is always yours:
      the provider runs the machines, but you remain accountable for what is in them and
      who may see it (\chref{ethics}).};

\node[anchor=west, align=left, text width=6.6cm, font=\scriptsize, text=gray!55!black]
     at (5.5,-2.5)
     {\textbf{Examples:} IaaS --- a GPU virtual machine you configure yourself. PaaS ---
      a managed notebook or training service. SaaS --- a BI dashboard or a hosted model
      API. Moving right buys convenience and costs control.};
\end{tikzpicture}}%
{The service models as a stack of responsibilities. Reading which layers are yours tells
you exactly what you must secure, patch and pay for.}{service-models}

\section{Virtual Machines, Containers, Serverless}

\begin{definitionbox}[Container]
A package containing an application together with its libraries and configuration,
sharing the host operating system kernel rather than running a full guest OS. Starts
in under a second, measured in megabytes, and runs identically wherever it is
deployed.
\end{definitionbox}

\dsfig{%
\begin{tikzpicture}[font=\scriptsize,
  l/.style={rounded corners=2pt, draw=#1, fill=#1!14, text=#1!50!black,
            minimum height=0.45cm, font=\tiny},
  hd/.style={font=\tiny\bfseries, text=primary}]

% --- VMs
\node[hd] at (2.0,3.3) {VIRTUAL MACHINES};
\foreach \i in {0,1,2}{
  \pgfmathsetmacro{\xx}{\i*1.35}
  \node[l=accent, minimum width=1.2cm] at (\xx,2.75) {app};
  \node[l=goldhl, minimum width=1.2cm] at (\xx,2.28) {libs};
  \node[l=secondary, minimum width=1.2cm, minimum height=0.6cm] at (\xx,1.75) {guest OS};}
\node[l=purplehl, minimum width=3.9cm] at (1.35,1.2) {hypervisor};
\node[l=primary, minimum width=3.9cm] at (1.35,0.73) {host OS};
\node[l=neutral, minimum width=3.9cm] at (1.35,0.26) {hardware};
\node[font=\tiny, text=accent!75!black, align=center, text width=4.0cm] at (1.35,-0.45)
     {each carries a whole OS:\\GBs, slow to start,\\strong isolation};

% --- Containers
\node[hd] at (8.4,3.3) {CONTAINERS};
\foreach \i in {0,...,4}{
  \pgfmathsetmacro{\xx}{6.4+\i*0.95}
  \node[l=accent, minimum width=0.85cm] at (\xx,2.75) {app};
  \node[l=goldhl, minimum width=0.85cm] at (\xx,2.28) {libs};}
\node[l=greenhl, minimum width=4.7cm] at (8.3,1.75) {container runtime};
\node[l=primary, minimum width=4.7cm] at (8.3,1.24) {host OS (shared kernel)};
\node[l=neutral, minimum width=4.7cm] at (8.3,0.73) {hardware};
\node[font=\tiny, text=greenhl!45!black, align=center, text width=4.6cm] at (8.3,-0.45)
     {share the kernel: MBs,\\start in $<1$s, many per host,\\weaker isolation};

\node[anchor=west, align=left, text width=13.6cm, font=\tiny, text=primary!60!black]
     at (-1.0,-1.6)
     {\textbf{Why this matters to a data scientist specifically:} ``it worked on my
      laptop'' is a reproducibility failure, and a container is the fix. The image pins
      the Python version, every library version and the system libraries beneath them, so
      the model that trained in March can be retrained identically in November. A
      \texttt{requirements.txt} does not achieve this; an image does.};
\end{tikzpicture}}%
{Virtual machines versus containers. For data science the decisive property is not
efficiency but reproducibility: the container image is an exact, versioned record of the
environment a result was produced in.}{vm-container}

\begin{conceptbox}[Serverless]
You supply a function; the provider runs it on demand and bills per invocation.
Nothing is running --- or being paid for --- between calls. Excellent for
intermittent, short, stateless work such as scoring a single record on arrival.
Poor for long training jobs, for anything needing a GPU, and for workloads
sensitive to the \term{cold start} delay on the first call after an idle period.
\end{conceptbox}

\section{Cloud, Fog and Edge}

\begin{definitionbox}[Edge, Fog and Cloud]
\term{Edge}: computation on or beside the device that generates the data.
\term{Fog}: an intermediate tier --- a gateway or local server serving a site.
\term{Cloud}: centralised, effectively unlimited compute, far away in network terms.
\end{definitionbox}

\dsfig{%
\begin{tikzpicture}[font=\scriptsize,
  tier/.style={rounded corners=4pt, draw=#1, fill=#1!8, text width=3.5cm, align=center,
               font=\tiny, inner sep=5pt, minimum height=2.85cm},
  hd/.style={rounded corners=3pt, fill=#1, text=white, font=\tiny\bfseries,
             text width=3.5cm, align=center, minimum height=0.55cm}]

\node[hd=greenhl] at (0,0)    {EDGE --- on the device};
\node[tier=greenhl] at (0,-1.8){%
  \textbf{Latency:} $<10$ ms\\
  \textbf{Compute:} tiny (MCU)\\
  \textbf{Cost:} no bandwidth\\[4pt]
  \emph{Anomaly threshold on a vibration sensor; wake-word detection}\\[4pt]
  {\color{greenhl!45!black}works when the network is down}};

\node[hd=goldhl] at (4.7,0)    {FOG --- site gateway};
\node[tier=goldhl] at (4.7,-1.8){%
  \textbf{Latency:} 10--50 ms\\
  \textbf{Compute:} moderate\\
  \textbf{Cost:} low bandwidth\\[4pt]
  \emph{Aggregate 400 sensors; run a small model for the whole factory}\\[4pt]
  {\color{goldhl!50!black}sees the whole site, not one device}};

\node[hd=secondary] at (9.4,0)    {CLOUD --- central};
\node[tier=secondary] at (9.4,-1.8){%
  \textbf{Latency:} 50--500 ms\\
  \textbf{Compute:} unlimited, GPUs\\
  \textbf{Cost:} bandwidth + compute\\[4pt]
  \emph{Train models on the whole fleet; long-term storage; dashboards}\\[4pt]
  {\color{secondary!55!black}the only place big training happens}};

\draw[-{Stealth[length=2.4mm]}, primary!35, line width=5pt, opacity=0.45]
     (-1.9,-3.65) -- (11.3,-3.65);
\node[font=\tiny\bfseries, text=primary!60!black] at (4.7,-4.0)
     {more compute, more latency, more bandwidth cost $\longrightarrow$};

\node[anchor=west, align=left, text width=13.5cm, font=\tiny, text=gray!55!black]
     at (-1.9,-4.75)
     {\textbf{The standard pattern is not to choose but to split:} \emph{train in the
      cloud, infer at the edge}. The heavy, occasional work goes where compute is cheap;
      the light, constant work goes where latency is low and bandwidth is expensive. This
      is the architecture of nearly every production IoT system (\chref{iot}).};
\end{tikzpicture}}%
{The three tiers. Placement is decided by the latency the decision requires and the
bandwidth you can afford --- not by where the computation is most convenient to
write.}{cloud-fog-edge}

\begin{worked}[Where Should the Model Run?]
400 motors, each producing 6 channels at 100 Hz. A vibration anomaly must trigger a
shutdown.

\smallskip
\textbf{Option A --- send everything to the cloud.}\\
Data rate: $400 \times 6 \times 100 \times 4\text{ bytes} = 960$ kB/s $\approx$
83 GB/day $\approx$ 30 TB/year. Round-trip latency: 100--300 ms. Also: if the link
drops, the factory has no protection at all.

\textbf{Option B --- everything at the edge.}\\
Latency: under 10 ms; no bandwidth cost; works offline. But a microcontroller
cannot train, and cannot see the other 399 motors, so it can never learn what
fleet-wide normal looks like.

\textbf{Option C --- split (the answer).}
\begin{tight}
  \item \emph{Edge:} compute window statistics on-device and run a quantised
        threshold model. Trigger shutdown locally in under 10 ms. Transmit only
        features and alerts --- roughly 1 kB per device per minute, about
        \textbf{0.6 GB/year} in total.
  \item \emph{Fog:} the site gateway correlates across all 400 motors, catching
        site-wide events no single device could see.
  \item \emph{Cloud:} retrain the model monthly on the accumulated feature history
        and push the updated model back to the devices.
\end{tight}

\smallskip
\textbf{Result:} bandwidth falls by a factor of roughly 50{,}000, the safety-critical
decision meets its latency budget, protection survives a network outage, and the
model still improves from fleet-wide learning. Note that the constraint that decided
this was \emph{latency and availability}, not accuracy.
\end{worked}

\section{Cost and the Shared Responsibility Model}

\begin{alertbox}[The Costs That Surprise Students]
\begin{tight}
  \item \textbf{Idle GPUs.} A GPU instance left running overnight costs the same as
        one training. Shut them down; use spot instances for interruptible work.
  \item \textbf{Egress.} Moving data \emph{into} a cloud is usually free; moving it
        \emph{out} is charged per gigabyte. This is what makes Option A above
        expensive forever, not just once.
  \item \textbf{Cross-region traffic.} Storage in one region and compute in another
        pays egress on every read.
  \item \textbf{Forgotten storage.} Nobody deletes the 2 TB of intermediate
        Parquet from March. Set lifecycle rules on creation.
\end{tight}
\end{alertbox}

\begin{definitionbox}[Shared Responsibility]
The provider secures the cloud --- physical sites, hardware, hypervisor, managed
services. You secure what is \emph{in} the cloud --- your data, access control,
network configuration and application code. Nearly every publicised ``cloud breach''
is a failure on the customer's side of this line, most often a misconfigured storage
bucket.
\end{definitionbox}

%---------------------------------------------------------------------
\begin{fourlenses}{Cloud and Edge}
\lensLA{Almost entirely cloud, since inference is nightly rather than real time and
volumes are modest. The dominant constraint is not technical but legal:
\textbf{data residency}. Student records may be required to remain in a specific
jurisdiction, which decides your region before any performance consideration.
Encrypt at rest and in transit, and restrict access by role --- most staff need
aggregates, not individual records.}

\lensPM{Nearly all the tooling is SaaS --- tracker, CI, chat --- so the data science
work is integration rather than infrastructure. The relevant risks are lock-in and
API stability: a vendor changing its API without notice breaks your pipeline
(\chref{dataeng}), and export capability should be checked before adoption, not
after.}

\lensIOT{This lens is the whole reason the chapter exists, and the worked example is
its canonical form: train in the cloud, infer at the edge, aggregate at the fog.
Devices must also be updatable over the air, because a model that cannot be
redeployed cannot respond to the drift of \chref{timeseries}. Budget bandwidth per
device before choosing a model size --- on cellular-connected fleets, bandwidth,
not accuracy, sets the architecture.}

\lensSEC{Cloud brings both new defences and a much larger attack surface. Log
collection benefits enormously from elastic storage, and managed detection services
are genuinely good. But the shared responsibility line is where organisations fail:
public buckets, over-permissive IAM roles, and unencrypted snapshots are all on the
customer's side. Add to that the \textbf{multi-tenancy} question --- your workload
shares hardware with strangers --- and containers' weaker isolation compared with
virtual machines becomes a real consideration for sensitive processing.}
\end{fourlenses}

%---------------------------------------------------------------------
\begin{lab}[Containerise a Model for Reproducibility]
\begin{lstlisting}[language=bash]
# ---- Dockerfile: an exact, versioned record of the environment --------
FROM python:3.11-slim                    # pin the version. never use :latest

WORKDIR /app
COPY requirements.txt .
RUN pip install --no-cache-dir -r requirements.txt

COPY model.pkl serve.py ./
EXPOSE 8080
CMD ["python", "serve.py"]
\end{lstlisting}

\begin{lstlisting}[language=Python]
# ---- serve.py: a minimal scoring service -----------------------------
import joblib, pandas as pd
from flask import Flask, request, jsonify

app = Flask(__name__)
model = joblib.load("model.pkl")

@app.route("/health")                     # every service needs this
def health():
    return {"status": "ok"}, 200

@app.route("/predict", methods=["POST"])
def predict():
    df = pd.DataFrame(request.get_json()["rows"])
    proba = model.predict_proba(df)[:, 1]
    return jsonify({"scores": proba.tolist()})

if __name__ == "__main__":
    app.run(host="0.0.0.0", port=8080)
\end{lstlisting}

\begin{lstlisting}[language=bash]
# ---- build, run, and shrink for the edge -----------------------------
docker build -t risk-model:2026.08.1 .    # tag with a version, not "latest"
docker run -p 8080:8080 risk-model:2026.08.1

# For an edge device, the same model must be made small:
#   quantise float32 -> int8   (roughly 4x smaller, 2-3x faster)
#   python -m tf.lite.TFLiteConverter ... with optimizations=[OPTIMIZE_FOR_SIZE]
# Then measure accuracy AFTER quantisation -- it usually drops a little,
# and whether that drop is acceptable is a Chapter 14 question.
\end{lstlisting}
\end{lab}

%---------------------------------------------------------------------
\begin{checkpoint}
\begin{tightnum}
  \item Under the shared responsibility model, who is accountable for a publicly
        readable storage bucket containing student records?
  \item Give two reasons why a container is a better answer to ``it worked on my
        laptop'' than a \texttt{requirements.txt} file.
  \item A safety shutdown must occur within 20 ms of detection. Which tier must the
        inference run on, and why is that not negotiable?
\end{tightnum}
\end{checkpoint}

%---------------------------------------------------------------------
\begin{chaptersummary}
\begin{tight}
  \item IaaS, PaaS and SaaS differ in how much of the stack you manage --- but the
        data layer is always yours.
  \item Containers share the host kernel, start in under a second, and pin the
        entire environment, which is what makes results reproducible.
  \item Serverless suits short, intermittent, stateless work; not training, not GPUs.
  \item Edge, fog and cloud trade latency against compute. The standard pattern is
        train in the cloud, infer at the edge.
  \item Placement is decided by latency requirements, bandwidth cost and offline
        availability --- not by convenience.
  \item Egress, idle GPUs and forgotten storage are the costs that surprise people;
        misconfiguration on the customer's side causes most cloud breaches.
\end{tight}
\end{chaptersummary}

\begin{reviewq}
  \item \textbf{(MCQ)} Containers differ from virtual machines principally because
        they: (a)~are more secure (b)~share the host kernel (c)~cannot run Linux
        (d)~require a GPU.
  \item \textbf{(MCQ)} Which workload is least suited to serverless? (a)~scoring one
        record on arrival (b)~a nightly 6-hour GPU training job (c)~resizing an
        uploaded image (d)~responding to a webhook.
  \item \textbf{(Short)} Explain data residency and give a scenario where it
        overrides a performance-based choice of region.
  \item \textbf{(Short)} Why is egress cost a recurring rather than a one-off
        concern in an IoT architecture?
  \item \textbf{(Short)} What is a cold start, and which deployment model suffers
        from it?
  \item \textbf{(Applied)} A hospital wants real-time patient-monitoring alerts.
        Specify which computation runs at the edge, fog and cloud, and justify each
        placement by latency, bandwidth, privacy and availability.
  \item \textbf{(Applied)} Estimate the annual bandwidth for 250 devices sending
        4 channels at 50 Hz as 4-byte floats. Then propose an edge-feature scheme
        and estimate the new figure, showing your working.
\end{reviewq}
